% ex: ts=2 sw=2 sts=2 et filetype=tex
% SPDX-License-Identifier: CC-BY-SA-4.0

\chapter{Marco Teórico}
\label{ch:marcoteorico}

%% Quita la siguiente linea es solo para hacer texto aleatorio.
\lipsum[1]
\cite{knuthwebsite}

\section{Sistema web}

\subsection{PHP (Hypertext Preprocessor)}

Para el desarrollo de la lógica del servidor y la gestión de la
persistencia de datos, se seleccionó PHP (Hypertext Preprocessor).
PHP es un lenguaje de programación de código abierto, interpretado y
de propósito general, especialmente diseñado para el desarrollo web
backend y la generación de contenido dinámico \parencite{php_group_2026}.

Desde su creación, PHP ha evolucionado de ser un conjunto de herramientas
básicas a convertirse en un lenguaje robusto y fuertemente tipado en sus
versiones más recientes (PHP 8.x), ofreciendo un alto rendimiento, mejoras
en el consumo de memoria y características avanzadas como la compilación
Just-In-Time (JIT).

La adopción de esta tecnología en el presente proyecto se justifica
principalmente por los siguientes factores:

\begin{itemize}
	\item Amplio ecosistema y soporte: Al ser uno de los lenguajes
		más utilizados en la infraestructura web global, cuenta con una
		comunidad madura que garantiza una extensa documentación, librerías
		de terceros y parches de seguridad constantes.

	\item Compatibilidad y portabilidad: PHP ofrece una integración nativa
		con una gran variedad de sistemas de gestión de bases de datos
		relacionales (como MySQL y PostgreSQL) y puede ejecutarse de manera
		eficiente en múltiples sistemas operativos (Linux, Windows, macOS)
		y servidores web (Apache, Nginx).

	\item Sinergia con frameworks modernos: Su arquitectura orientada a
		objetos proporciona la base estructural necesaria para soportar
		entornos de desarrollo de alto nivel, como el framework Laravel,
		optimizando los tiempos de respuesta del servidor y facilitando
		el mantenimiento del código a largo plazo.
\end{itemize}

\subsection{El Framework Laravel}

Para la estructuración y el desarrollo de la aplicación web, se
implementó Laravel, un framework de código abierto basado en PHP y
diseñado bajo el patrón de arquitectura Modelo-Vista-Controlador (MVC)
\parencite{otwell_laravel_2026}. Desde su concepción, Laravel ha tenido
como objetivo simplificar las tareas comunes en el desarrollo de software
web —como la autenticación, el enrutamiento, las sesiones y el
almacenamiento en caché— sin sacrificar la flexibilidad ni la
escalabilidad del sistema.

La integración de Laravel en la arquitectura del proyecto aporta ventajas
metodológicas y técnicas fundamentales:

\begin{itemize}
	\item Abstracción de la Base de Datos: A través de Eloquent ORM
		(Object-Relational Mapping), el framework permite interactuar
		con la base de datos mediante programación orientada a objetos
		en lugar de consultas SQL nativas. Esto no solo incrementa la
		seguridad contra ataques de inyección SQL, sino que además
		facilita la migración y consistencia de los datos mediante el
		uso de Migrations.

	\item Seguridad Integrada: Laravel provee mecanismos robustos
		y preconfigurados para mitigar las vulnerabilidades web más
		comunes, tales como la falsificación de peticiones en sitios
		cruzados (CSRF), la secuencia de comandos en sitios cruzados
		(XSS) y la encriptación de datos sensibles.

	\item Ecosistema y Modularidad: La herramienta cuenta con un
		sistema de empaquetado y gestión de dependencias (a través de
		Composer) que permite la reutilización de componentes y la
		integración de servicios externos de manera limpia, acelerando
		el ciclo de desarrollo y garantizando un código mantenible.
\end{itemize}

Al combinar la robustez y velocidad de PHP en su versión actual
con la estructura organizada de Laravel, se garantiza una plataforma
backend eficiente, segura y adaptable a futuras expansiones o
requerimientos del sistema.

\section{Plataforma de hardware utilizada}

Para el despliegue físico y la ejecución del sistema, se seleccionó la
computadora de placa única (SBC, por sus siglas en inglés Single Board
Computer) Raspberry Pi. Este hardware de código abierto está diseñado
sobre una arquitectura ARM y destaca por su tamaño compacto, bajo
consumo energético y su capacidad para ejecutar sistemas operativos
basados en Linux de manera nativa \parencite{raspberry_pi_2026}.

\subsection{Raspberry Pi como Plataforma de Despliegue}

En el contexto de este proyecto, la Raspberry Pi actúa como el servidor
local (o nodo central), encargándose de alojar el servidor web, el
intérprete de PHP y el motor de bases de datos. Su inclusión en la
arquitectura de hardware se justifica por los siguientes aspectos
técnicos y metodológicos:

\begin{itemize}
	\item Eficiencia de Recursos y Costo: Ofrece una relación
		costo-rendimiento óptima para proyectos de investigación y prototipado.
		Permite simular un entorno de producción real en un servidor dedicado
		sin incurrir en los costos recurrentes de infraestructura en la nube.

	\item Compatibilidad con el Entorno LAMP/LEMP: Al ejecutar distribuciones
		oficiales como Raspberry Pi OS (basada en Debian), el hardware es
		plenamente compatible con servidores web como Apache o Nginx,
		permitiendo el montaje del entorno necesario para ejecutar Laravel de
		forma idéntica a un servidor web comercial.

	\item Consumo Energético y Portabilidad: Su bajo requerimiento de
		energía (típicamente alimentado por USB-C a 5V) y su reducido factor
		de forma facilitan el despliegue del sistema en entornos aislados o
		embebidos, lo que aporta un valor agregado de portabilidad a la
		solución propuesta.
\end{itemize}

Con la integración de este hardware, el proyecto demuestra que la combinación
de PHP y Laravel no solo es viable en servidores empresariales, sino que
también es capaz de operar de manera fluida y eficiente en arquitecturas
de hardware limitado, optimizando el uso de la memoria RAM y el
procesamiento del procesador ARM.


